\documentclass[a4paper,12pt]{article}
\usepackage[T2A]{fontenc}
\usepackage[utf8]{inputenc}
\usepackage[russian]{babel}
\usepackage{amsmath,amssymb,mathtools}
\usepackage{geometry}
\geometry{left=24mm,right=24mm,top=24mm,bottom=25mm}
\usepackage{microtype}
\usepackage{booktabs,tabularx,longtable,array}
\usepackage{enumitem}
\usepackage{tikz}
\usetikzlibrary{arrows.meta,positioning,shapes.geometric,calc}
\usepackage{hyperref}
\usepackage{fancyhdr}
\usepackage{xcolor}
\usepackage{parskip}
\definecolor{accent}{HTML}{008080}
\definecolor{darkaccent}{HTML}{005354}
\definecolor{quiet}{HTML}{51585A}
\hypersetup{colorlinks=true,linkcolor=darkaccent,urlcolor=darkaccent}
\setlength{\parindent}{0pt}
\setlength{\parskip}{5pt plus 1pt minus 1pt}
\setlist[itemize]{leftmargin=1.4em,itemsep=2pt,topsep=3pt}
\setlist[enumerate]{leftmargin=1.6em,itemsep=9pt,topsep=7pt}
\pagestyle{fancy}
\fancyhf{}
\fancyhead[L]{\small\color{quiet} Степени и физические формулы}
\fancyhead[R]{\small\color{quiet} Даниил Хрястов}
\fancyfoot[C]{\small\color{quiet}\thepage}
\renewcommand{\headrulewidth}{0.35pt}
\setlength{\headheight}{15pt}
\newcommand{\sect}[1]{\section{#1}\vspace{-0.35em}\noindent{\color{accent}\rule{\linewidth}{0.45pt}}\vspace{0.25em}}
\newcommand{\key}[1]{{\color{darkaccent}\textbf{#1}}}
\newcommand{\ex}[1]{\par\medskip\noindent{\color{darkaccent}\textbf{#1}}\par\smallskip}
\renewcommand{\arraystretch}{1.25}
\begin{document}
\begin{titlepage}
\thispagestyle{empty}
\begin{center}
\vspace*{28mm}
{\large\color{darkaccent} МАТЕМАТИКА}\par
\vspace{25mm}
{\Huge Чек-лист}\par
\vspace{17mm}
{\LARGE\color{accent} Степенные преобразования\\[3pt] в задачах с физическими формулами}\par
\vspace{18mm}
{\large Прикладные задачи: степени, формулы, закономерности}\par
\vfill
{\Large Даниил Хрястов}\par
\vspace{12mm}
{\large 2 ноября 2026 г.}\par
\vspace{22mm}
{\normalsize Лёвин Артём Александрович\\[3pt] эксклюзивно для Даниила Хрястова}\par
\vspace{21mm}
\begin{tikzpicture}[x=1cm,y=1cm]
\draw[accent, line width=0.75pt] (0,0) .. controls (1.5,0.52) and (2.75,-0.39) .. (4.0,0)
.. controls (5.5,0.42) and (7.1,-0.35) .. (8.8,0);
\fill[darkaccent] (4.0,0) circle (1.25pt);
\end{tikzpicture}
\end{center}
\end{titlepage}

\sect{Главная идея занятия}
Во многих прикладных задачах нужная величина скрыта в степени. Знакомые действия --- раскрыть скобки, перенести слагаемое, разделить на коэффициент --- помогают \emph{выделить эту степень}. Затем используются свойства показателей и обратное возведение в степень.

\key{Маршрут решения:} запишите формулу $\longrightarrow$ подставьте данные $\longrightarrow$ выделите степень $\longrightarrow$ приведите числа к удобному виду $\longrightarrow$ найдите величину $\longrightarrow$ проверьте результат.

\sect{Свойства степеней, которые нужны в задачах}
Для $a>0$ и допустимых показателей:
\[
\begin{aligned}
a^m\cdot a^n&=a^{m+n}, &\frac{a^m}{a^n}&=a^{m-n},\\[2pt]
(a^m)^n&=a^{mn}, &(ab)^n&=a^n b^n,\\[2pt]
a^{-n}&=\frac1{a^n}, &a^{\frac mn}&=\sqrt[n]{a^m}.
\end{aligned}
\]
При решении физических задач рассматривайте положительные значения масс, объёмов, температур и других величин, если это следует из смысла условия.

\ex{Пример 1. Как убрать дробный показатель}
Если $x>0$ и $x^{5/3}=32$, нужно получить степень $x^1$:
\[
 x=32^{3/5}=(2^5)^{3/5}=2^3=8.
\]
Возведение обеих частей в степень $3/5$ допустимо, поскольку $x>0$.

\ex{Пример 2. Отрицательная степень}
\[
 \frac{1}{125}=\frac{1}{5^3}=5^{-3},\qquad
 (5^{-5})^{3/5}=5^{-3}=\frac1{125}.
\]
\textbf{Проверка:} показатель степени умножается на $3/5$, поэтому $-5\cdot\frac35=-3$.

\sect{Стандартный вид числа и порядок действий}
Большие и малые числа удобно записывать как $a\cdot10^n$, где $1\leqslant |a|<10$:
\[
 0{,}000000057=5{,}7\cdot10^{-8}=57\cdot10^{-9},
 \qquad 625000=6{,}25\cdot10^5.
\]
При умножении показатели десяток складываются, при делении вычитаются:
\[
 \frac{10^{21}}{10^{-9}\cdot10^{20}}
 =10^{21-(-9)-20}=10^{10}.
\]
\key{Важно:} при переносе дроби через знак равенства не «переворачивайте всё выражение» механически. Из $A=B/C$ следует $AC=B$ при $C\ne0$.

\sect{Куб неизвестной длины}
На занятии встречалась формула вида
\[
 F=\rho g\ell^3-mg.
\]
Здесь $F$ --- сила, $\rho$ --- плотность, $g$ --- ускорение свободного падения, $m$ --- масса, $\ell$ --- искомая длина.

\ex{Пример 3. Выделяем $\ell^3$}
Пусть $F=1000\,\text{Н}$, $\rho=1000\,\text{кг/м}^3$, $g=10\,\text{м/с}^2$, $m=25\,\text{кг}$. Тогда
\[
\begin{aligned}
1000&=1000\cdot10\cdot\ell^3-25\cdot10,\\
1000+250&=10000\ell^3,\\
\ell^3&=\frac{1250}{10000}=\frac18,\\
\ell&=\sqrt[3]{\frac18}=\frac12=0{,}5\ \text{м}.
\end{aligned}
\]
\textbf{Проверка:} $1000\cdot10\cdot(0{,}5)^3-25\cdot10=1250-250=1000$ Н.

\sect{Четвёртая степень в формуле излучения}
В модели $P=\sigma ST^4$ величина $T$ скрыта в четвёртой степени. Если $\sigma>0$, $S>0$ и $P>0$, то
\[
 T^4=\frac{P}{\sigma S},\qquad
 \boxed{T=\left(\frac{P}{\sigma S}\right)^{1/4}}.
\]
В физическом смысле выбирается положительное значение температуры.

\ex{Пример 4. Большие и малые степени десяти}
По данным, разобранным на доске:
\[
 P=15625\cdot10^{21},\quad
 \sigma=57\cdot10^{-9},\quad S=\frac{10^{20}}{228}.
\]
\[
\begin{aligned}
T^4&=\frac{15625\cdot10^{21}\cdot228}{57\cdot10^{-9}\cdot10^{20}}
 =15625\cdot4\cdot10^{10}\\
 &=5^6\cdot2^2\cdot(2\cdot5)^{10}
 =5^{16}\cdot2^{12},\\
T&=(5^{16}\cdot2^{12})^{1/4}
 =5^4\cdot2^3=5000.
\end{aligned}
\]
\key{Приём:} если нужно извлечь корень четвёртой степени, старайтесь получить степени с показателями, кратными четырём.

\sect{Закон радиоактивного распада}
Для массы изотопа применялась модель
\[
 \boxed{m(t)=m_0\cdot2^{-t/\tau}},
\]
где $m_0>0$ --- начальная масса, $\tau>0$ --- период полураспада, $t\geqslant0$ --- время.

\ex{Пример 5. Находим время}
Пусть $m_0=40$ мг, $\tau=10$ мин. Когда останется $5$ мг?
\[
 5=40\cdot2^{-t/10}
 \quad\Longrightarrow\quad
 2^{-t/10}=\frac5{40}=\frac18=2^{-3}.
\]
Основания равны и отличны от единицы, поэтому показатели равны:
\[
 -\frac{t}{10}=-3 \quad\Longrightarrow\quad
 \boxed{t=30\ \text{мин}}.
\]
\textbf{Проверка:} три периода полураспада дают $40\to20\to10\to5$ мг.

\sect{Степень объёма в законе для газа}
На доске применялось соотношение
\[
 \boxed{pV^k=C},\qquad p>0,\quad V>0,\quad C>0.
\]
Если $k\ne0$ известно, объём находится так:
\[
 V^k=\frac Cp,\qquad V=\left(\frac Cp\right)^{1/k}.
\]
\ex{Пример 6. Дробная степень объёма}
Даны $k=5/3$, $p=6{,}25\cdot10^5$ Па, $C=200$:
\[
\begin{aligned}
V^{5/3}&=\frac{200}{625000}=\frac1{3125}=5^{-5},\\
V&=(5^{-5})^{3/5}=5^{-3}=\frac1{125}
=\boxed{0{,}008\ \text{м}^3}.
\end{aligned}
\]
\key{Контроль:} $V$ положителен, и $625000\cdot(1/125)^{5/3}=200$.

\sect{Сравнение двух состояний газа}
Если для газа постоянна величина $pV^k$, то
\[
 p_1V_1^k=p_2V_2^k.
\]
\ex{Пример 7. Находим неизвестный показатель}
При $V_1=2V_2$ и $p_2=4p_1$ получаем
\[
 p_1(2V_2)^k=4p_1V_2^k.
\]
Поскольку $p_1>0$ и $V_2>0$, разрешено сократить одинаковые множители:
\[
 2^k=4=2^2\quad\Longrightarrow\quad\boxed{k=2}.
\]
\key{Замечание:} равенство $2^k=2^2$ позволяет сравнить показатели. При разных основаниях такой переход был бы неверен.

\sect{Типичные ошибки и самопроверка}
\begin{itemize}
\item $\displaystyle\frac{a}{b/c}=\frac{ac}{b}$ при $b\ne0$, $c\ne0$: важно учитывать весь знаменатель.
\item $10^{-9}$ при делении даёт вычитание отрицательной девятки: $-(-9)=+9$.
\item Из $V^{5/3}=q>0$ следует $V=q^{3/5}$, а не $q^{5/3}$.
\item При $a^r=a^s$ сравнивайте показатели лишь при $a>0$, $a\ne1$.
\item Не сокращайте $p_1$ или $V_2^k$, не указав, почему они не равны нулю.
\item После решения проверьте исходную формулу, знак и единицы измерения.
\end{itemize}

\clearpage
\thispagestyle{fancy}
\begin{center}
{\Large\color{darkaccent} Алгоритм: как найти неизвестную в степенной формуле}\par
\vspace{12mm}
\begin{tikzpicture}[
  node distance=9mm and 8mm,
  arr/.style={-{Latex[length=2.7mm]},line width=.8pt,draw=darkaccent},
  act/.style={rectangle,draw=accent,line width=.7pt,rounded corners=1pt,align=center,text width=9.5cm,minimum height=12mm,inner sep=5pt},
  side/.style={rectangle,draw=accent,line width=.7pt,rounded corners=1pt,align=center,text width=5.6cm,minimum height=20mm,inner sep=5pt},
  decision/.style={diamond,draw=darkaccent,line width=.7pt,aspect=2.7,align=center,text width=4.6cm,inner sep=3pt},
  start/.style={ellipse,draw=darkaccent,line width=.8pt,align=center,text width=5.8cm,minimum height=11mm}]
\node[start] (start) {Запишите формулу и данные};
\node[act,below=of start] (sub) {Подставьте числа, приведите их к степеням и выделите выражение с неизвестной};
\node[decision,below=11mm of sub] (choice) {Где неизвестная?};
\node[side,below left=14mm and 1mm of choice] (base) {\textbf{В основании степени}\\$x^r=A$\\Возведите обе части в степень $1/r$; учтите ограничения};
\node[side,below right=14mm and 1mm of choice] (expo) {\textbf{В показателе}\\$a^{f(x)}=a^c$\\Приведите к одному основанию $a>0$, $a\ne1$};
\node[act,below=22mm of choice,yshift=-38mm] (verify) {Выполните вычисления; проверьте подстановкой, знак и единицы};
\node[start,below=of verify] (end) {Запишите ответ};
\draw[arr] (start)--(sub);
\draw[arr] (sub)--(choice);
\draw[arr] (choice.south west) -| node[pos=.32,above left,font=\small]{основание} (base.north);
\draw[arr] (choice.south east) -| node[pos=.32,above right,font=\small]{показатель} (expo.north);
\draw[arr] (base.south) |- (verify.west);
\draw[arr] (expo.south) |- (verify.east);
\draw[arr] (verify)--(end);
\end{tikzpicture}
\end{center}
\clearpage

\sect{Тренировочный блок --- Путь А}
Выполните \textbf{10 заданий}, располагая решения последовательно. Для величин с физическим смыслом укажите единицы измерения. Сложность постепенно возрастает.

\begin{enumerate}[label=\textbf{\arabic*.}]
\item Упростите: $\displaystyle \frac{2^5\cdot2^{-3}}{2^{-2}}$.
\item Вычислите: $\displaystyle (4\cdot10^{-7})\cdot(2{,}5\cdot10^9)$.
\item Найдите положительное $x$, если $\displaystyle x^{3/2}=125$.
\item В формуле $F=\rho g\ell^3-mg$ известны $\rho=1000$ кг/м$^3$, $g=10$ м/с$^2$, $m=16$ кг, $F=2000$ Н. Найдите $\ell$ в метрах.
\item Масса изотопа меняется по закону $m(t)=m_0\cdot2^{-t/\tau}$. Найдите массу через $21$ минуту, если $m_0=48$ мг и $\tau=7$ мин.
\item Найдите $T>0$ из формулы $P=\sigma ST^4$, если $P=8\cdot10^6$, $\sigma=5\cdot10^{-8}$ и $S=10$ (в согласованных единицах).
\item Через сколько минут останется $5$ мг изотопа, если $m_0=80$ мг, $\tau=12$ мин и $m(t)=m_0\cdot2^{-t/\tau}$?
\item Для газа $pV^{5/3}=80$, а давление $p=2{,}5\cdot10^5$ Па. Найдите $V>0$ в м$^3$.
\item В двух состояниях газа $p_1V_1^k=p_2V_2^k$. Известно, что $V_1=3V_2$ и $p_2=9p_1$. Найдите $k$.
\item В газовом процессе $p_1V_1^{5/3}=p_2V_2^{5/3}$. Если $p_2=32p_1$, $V_1=0{,}016$ м$^3$, найдите $V_2$.
\end{enumerate}

\clearpage
\sect{Ответы к тренировочному блоку}
\begin{center}
\begin{tabularx}{0.88\linewidth}{@{}>{\bfseries}cX@{}}
\toprule
№ & Ответ \\
\midrule
1 & $16$ \\
2 & $1000$ \\
3 & $x=25$ \\
4 & $\ell=0{,}6$ м \\
5 & $6$ мг \\
6 & $T=2000$ \\
7 & $48$ мин \\
8 & $V=0{,}008$ м$^3$ \\
9 & $k=2$ \\
10 & $V_2=0{,}002$ м$^3$ \\
\bottomrule
\end{tabularx}
\end{center}

\vspace{12mm}
\sect{Краткая памятка}
\begin{tabularx}{\linewidth}{@{}>{\raggedright\arraybackslash}p{.36\linewidth}X@{}}
\toprule
Сигнал в условии & Полезное действие \\
\midrule
Большие числа, очень маленькие коэффициенты & Запишите через степени десяти и соберите показатели. \\
Величина под степенью & Сначала выделите степень; затем возведите в обратную. \\
Время в показателе $2^{-t/\tau}$ & Представьте отношение масс как степень двойки. \\
Сравниваются два состояния газа & Приравняйте $p_1V_1^k$ и $p_2V_2^k$, затем сократите положительные множители. \\
Получен физический ответ & Проверьте размерность, допустимость и исходное равенство. \\
\bottomrule
\end{tabularx}

\vfill
\begin{center}
{\small\color{quiet} Лёвин Артём Александрович эксклюзивно для Даниила Хрястова}
\end{center}
\end{document}
